% Main Overleaf document. Compiler: pdfLaTeX.
% Options: blue (default), cyan, dark, handout.
\documentclass[10pt,aspectratio=169]{dgist_beamer}
\usepackage{pgfplots}
\pgfplotsset{compat=1.18}

\title{Research Presentation}
\subtitle{Robotics and Mechatronics Engineering}
\author{Donghae Kim}
\institute{Control Group for Intelligent Multi-Agent Systems\\
  Daegu Gyeongbuk Institute of Science and Technology}
\date{Seminar / Conference\\\today}
% Footer text is empty by default. Uncomment either line to add your own:
\setfooterleft{Donghae Kim, title of the presentation}
\setfooterright{Conference name, date, location, ...}
% Optional: \setbeameroption{show notes}
% Optional: \AtBeginSection[]{\dgistsectionframe}

\begin{document}
\dgisttitleframe

\begin{frame}
  \frametitle{Outline}
  \tableofcontents
\end{frame}

\section{Research content}
\dgistsectionframe

\begin{frame}
  \frametitle{Problem formulation}
  \framesubtitle{An example layout for equations and assumptions}
  Consider a network of $N$ agents with local objectives $f_i$:
  \begin{equation*}
    \underset{x\in\mathcal{X}}{\operatorname{minimize}}
    \quad F(x)=\sum_{i=1}^{N} f_i(x).
  \end{equation*}
  \vspace{.12cm}
  \begin{itemize}
    \setlength{\itemsep}{.18cm}
    \item Each agent has access to its own objective and local measurements.
    \item Neighboring agents exchange information over a communication graph.
    \item The goal is to reach a common feasible solution.
  \end{itemize}
  \vspace{.20cm}
  \begin{block}{Research question}
    Which assumptions guarantee convergence of the distributed algorithm?
  \end{block}
\end{frame}

\begin{frame}
  \frametitle{Assumptions and analysis}
  \begin{columns}[T,onlytextwidth]
    \begin{column}{.47\textwidth}
      \begin{block}{Assumptions}
        \begin{itemize}
          \setlength{\itemsep}{.15cm}
          \item Convex local objectives
          \item A nonempty feasible set
          \item Suitable graph connectivity
        \end{itemize}
      \end{block}
      \vspace{.18cm}
      State the exact assumptions for your model and algorithm here.
    \end{column}
    \begin{column}{.47\textwidth}
      \begin{exampleblock}{Analysis outline}
        \begin{enumerate}
          \setlength{\itemsep}{.15cm}
          \item Establish boundedness.
          \item Bound disagreement.
          \item Show optimality.
        \end{enumerate}
      \end{exampleblock}
      \vspace{.18cm}
      \begin{alertblock}{Scope}
        This slide illustrates formatting, not a complete convergence result.
      \end{alertblock}
    \end{column}
  \end{columns}
\end{frame}

\begin{frame}
  \frametitle{Illustrative convergence curves}
  \framesubtitle{Synthetic functions for demonstrating the plot style}
  \begin{columns}[c,onlytextwidth]
    \begin{column}{.64\textwidth}
      \centering
      \begin{tikzpicture}
        \begin{axis}[
          width=\linewidth,height=.49\paperheight,
          xmin=0,xmax=10,ymin=0,ymax=1.05,
          xlabel={Time $t$},ylabel={Normalized error},
          tick label style={font=\scriptsize,color=text},
          label style={font=\small,color=text},
          axis line style={draw=muted},tick style={draw=muted},
          grid=major,grid style={draw=rulecolor},
          legend style={font=\scriptsize,draw=none,fill=pagebackground,text=text},
          legend pos=north east,
          samples=101,domain=0:10]
          \addplot[highlight,very thick] {exp(-.65*x)};
          \addlegendentry{$e^{-0.65t}$}
          \addplot[dgist_cyan,very thick,dashed] {exp(-.30*x)};
          \addlegendentry{$e^{-0.30t}$}
        \end{axis}
      \end{tikzpicture}
    \end{column}
    \begin{column}{.31\textwidth}
      \textbf{Figure layout}\par\medskip
      Use consistent units and distinguish curves by line style as well as color.
      \par\bigskip
      Replace this plot with your simulation or experimental results.
    \end{column}
  \end{columns}
\end{frame}

\begin{frame}
  \frametitle{Comparison table}
  \framesubtitle{A reusable format for assumptions and guarantees}
  \renewcommand{\arraystretch}{1.55}
  \begin{center}
    \begin{tabular}{@{}p{.25\textwidth}p{.29\textwidth}p{.29\textwidth}@{}}
      \toprule
      \textcolor{highlight}{\textbf{Property}} &
      \textcolor{highlight}{\textbf{Baseline}} &
      \textcolor{highlight}{\textbf{Proposed method}}\\
      \midrule
      Graph model & State the assumption & State the assumption\\
      Objective class & State the assumption & State the assumption\\
      Constraints & State the admissible set & State the admissible set\\
      Guarantee & State the result & State the result\\
      \bottomrule
    \end{tabular}
  \end{center}
  \vspace{.25cm}
  Include the relevant references and conditions when comparing methods.
\end{frame}

\section{Template reference}
\dgistsectionframe

\begin{frame}
  \frametitle{DGIST color palette}
  \framesubtitle{Official CMYK definitions from DGIST Branding \& Identity}
  \renewcommand{\arraystretch}{1.3}
  \begin{center}\small
    \begin{tabular}{@{}llcl@{}}
      \toprule
      Color name & Pantone reference & Swatch & CMYK (\%)\\
      \midrule
      \texttt{dgist\_cyan} & 2985 CVC & \textcolor{dgist_cyan}{\rule{1.1cm}{.24cm}} & 72 / 0 / 0 / 0\\
      \texttt{dgist\_blue} & 285 CVC & \textcolor{dgist_blue}{\rule{1.1cm}{.24cm}} & 91 / 43 / 0 / 0\\
      \texttt{dgist\_navy} & 280 CVC & \textcolor{dgist_navy}{\rule{1.1cm}{.24cm}} & 100 / 72 / 0 / 18\\
      \texttt{dgist\_gray} & Cool Grey 9 CVC & \textcolor{dgist_gray}{\rule{1.1cm}{.24cm}} & 0 / 0 / 0 / 65\\
      \texttt{dgist\_lightgray} & Cool Grey 2 CVC & \textcolor{dgist_lightgray}{\rule{1.1cm}{.24cm}} & 0 / 0 / 0 / 11\\
      \bottomrule
    \end{tabular}
  \end{center}
  \vspace{.15cm}
  {\small Use \texttt{highlight}, \texttt{text}, and \texttt{muted} for content
  that adapts to light and dark themes.}\par\medskip
  {\scriptsize\href{https://www.dgist.ac.kr/eng/sub01_02_04_01.do}{Source: DGIST official color system}}
  \note{Color specifications: https://www.dgist.ac.kr/eng/sub01_02_04_01.do.
    The source logo files are unchanged. The title and closing slides use
    the original logo arrangement at reduced size, slightly lower, with centered text.
    Department website: https://www.dgist.ac.kr/robot/.}
\end{frame}

\begin{frame}[fragile]
  \frametitle{Settings for your next presentation}
  \begin{block}{Document class}
\begin{verbatim}
\documentclass[10pt,aspectratio=169]{dgist_beamer}
\end{verbatim}
  \end{block}
  \begin{itemize}
    \setlength{\itemsep}{.18cm}
    \item Add \texttt{dark} for a dark background or \texttt{cyan} for a cyan accent.
    \item Add \texttt{handout} to collapse overlays for distribution.
    \item Footer text is empty by default. Frame numbers remain visible.
    \item Use \verb|\setfooterleft| and \verb|\setfooterright| to add footer text.
  \end{itemize}
  \vspace{.18cm}
  The \texttt{README.md} describes the available commands.
\end{frame}

\dgistclosingframe{CiMAS\\Robotics and Mechatronics Engineering, DGIST}
\end{document}
